\documentclass[11pt]{article}
\usepackage[margin=1in]{geometry}
\usepackage{amsmath,amsthm,amssymb,mathtools}
\usepackage[T1]{fontenc}
\usepackage[utf8]{inputenc}

\newtheorem{theorem}{Theorem}
\newtheorem{proposition}{Proposition}
\newtheorem{lemma}{Lemma}
\theoremstyle{remark}
\newtheorem{remark}{Remark}

\newcommand{\stepd}[1]{\partial_{#1}}
\newcommand{\sig}{\sigma}

\title{d5\_220: Surfaced-bundle closure of the \(\sigma_{23}\) frontier}
\author{}
\date{2026-03-21}

\begin{document}
\maketitle

\paragraph{Purpose.}
The uploaded bundle contains the previously missing source file
\texttt{scripts/torus\_nd\_d5\_endpoint\_latin\_repair.py} together with the common mode-switch module
\texttt{torus\_nd\_d5\_layer3\_mode\_switch\_common.py}.  Once these are surfaced, the
\(\sigma_{21}\to\sigma_{23}\) frontier no longer requires a separate locality conjecture: the relevant
second-step comparison can be read off directly from the builder.

\paragraph{Scope.}
This note proves the exact theorem implemented by the surfaced bundle builder for the three row prototypes
\[
\sigma_{21}=[4,1,4]/[1,2,2],\qquad
\sigma_{22}=[4,1,4]/[2,1,2],\qquad
\sigma_{23}=[4,1,4]/[2,2,1],
\]
with the standard choice \(w_0=s_0=0\).  The conclusion is stronger than the odd-only target: it holds for every
modulus \(m>2\).

\paragraph{Bundle semantics used.}
Let \(m>2\).  In the surfaced common module, \texttt{precompute\_m} builds \texttt{step\_by\_dir[d]} by adding
\(+1\) in coordinate direction \(d\) modulo \(m\).  We denote this map by \(\stepd{d}\).
Thus the five coordinate steps commute:
\[
\stepd{i}\stepd{j}=\stepd{j}\stepd{i}\qquad (0\le i,j\le 4).
\]
In the surfaced endpoint builder, \texttt{\_prepare\_m} creates the layer-1 source families and \texttt{\_build\_candidate}
creates the labeled rows by repeated application of \(\stepd{d}\) according to the chosen left/right words.

\begin{lemma}[Exact right-source family]
Fix color \(0\), \(w_0=0\), and \(s_0=0\).  The right \(R1\) source family used by all three prototypes is exactly
\[
S_R=\{x:\ \mathrm{layer}(x)=1,\ q(x)=m-1,\ w(x)=0,\ s(x)\neq 0\}.
\]
Equivalently, since \(s=w+u\) and \(w=0\),
\[
S_R=\{x:\ \mathrm{layer}(x)=1,\ q(x)=m-1,\ w(x)=0,\ u(x)\neq 0\}.
\]
Hence
\[
|S_R|=m(m-1).
\]
\end{lemma}

\begin{proof}
This is exactly the definition of \texttt{layer1\_sets[(w,s0)][0]} in \texttt{\_prepare\_m}: layer \(1\),
\(q=m-1\), prescribed \(w\), and exclusion of the single value \(s=s_0\).  For \((w_0,s_0)=(0,0)\) the right family
uses \(w=0\), so \(s=u\).  The free coordinates are \(v\in\mathbb Z/m\mathbb Z\) and \(u\in(\mathbb Z/m\mathbb Z)\setminus\{0\}\),
which gives \(m(m-1)\) states.
\end{proof}

\begin{lemma}[No overlap conflicts for \(\sigma_{21},\sigma_{22},\sigma_{23}\)]
For every \(m>2\), the three prototypes \(\sigma_{21},\sigma_{22},\sigma_{23}\) are built successfully by
\texttt{\_build\_candidate}; in particular, the layer-2 and layer-3 overlap tests vanish.
\end{lemma}

\begin{proof}
The left word is fixed as \([4,1,4]\), so for color \(0\) one has
\[
L2=\stepd{4}(S_L),\qquad L3=\stepd{1}(L2),
\]
where the left source family satisfies \(q=m-1\) and \(w=m-1\).  Therefore every point of \(L2\) has
\((q,w)=(m-1,m-1)\), while every point of \(L3\) has \((q,w)=(0,m-1)\).

For the right side one has:
\begin{align*}
\sigma_{21}:&\quad R2=\stepd{1}(S_R), && R3=\stepd{2}(R2),\\
\sigma_{22}:&\quad R2=\stepd{2}(S_R), && R3=\stepd{1}(R2),\\
\sigma_{23}:&\quad R2=\stepd{2}(S_R), && R3=\stepd{2}(R2).
\end{align*}
Since every source in \(S_R\) has \((q,w)=(m-1,0)\), this gives:
\begin{align*}
\sigma_{21}:&\quad R2:(q,w)=(0,0), && R3:(q,w)=(0,1),\\
\sigma_{22}:&\quad R2:(q,w)=(m-1,1), && R3:(q,w)=(0,1),\\
\sigma_{23}:&\quad R2:(q,w)=(m-1,1), && R3:(q,w)=(m-1,2).
\end{align*}
Thus \(L2\cap R2=\varnothing\) for all \(m>1\), and \(L3\cap R3=\varnothing\) for all \(m>2\).
Therefore the builder never triggers its layer-2 or layer-3 conflict aborts on these three prototypes.
\end{proof}

\begin{lemma}[The rooted local pattern is constant]
For every source \(x\in S_R\), the outgoing neighbor-label pattern of \(\sigma_{21}\) in directions
\(0,1,2,3,4\) is exactly
\[
[B,R2,B,B,B].
\]
\end{lemma}

\begin{proof}
By construction, \(\stepd{1}(x)\in R2\), so direction \(1\) has label \(R2\).
All other outgoing neighbors \(\stepd{d}(x)\) lie in layer \(2\), hence cannot belong to \(L1\) or \(R1\).
They also fail the explicit \((q,w)\)-characterizations of \(L2\) and \(R2\):
\begin{align*}
\stepd{0}(x)&:(q,w)=(m-1,0),\\
\stepd{2}(x)&:(q,w)=(m-1,1),\\
\stepd{3}(x)&:(q,w)=(m-1,0),\\
\stepd{4}(x)&:(q,w)=(m-1,0),
\end{align*}
whereas \(L2\) has \((q,w)=(m-1,m-1)\) and \(R2\) has \((q,w)=(0,0)\) for \(\sigma_{21}\).
They cannot lie in \(L3\) or \(R3\) either because those are layer-3 sets.  Hence only direction \(1\) is occupied,
as claimed.
\end{proof}

\begin{theorem}[Surfaced-bundle closure of the \(\sigma_{23}\) frontier]
For every modulus \(m>2\), every \(\sigma_{21}\) immediate \(R1\) source satisfies
\[
\sigma_{21}^2(x)=\sigma_{22}^2(x)\neq \sigma_{23}^2(x).
\]
Equivalently, the actual second step of the \(\sigma_{21}\) immediate \(R1\) family always realizes the
\(\sigma_{22}\) prototype and never the \(\sigma_{23}\) prototype.  In particular, the odd-modulus
\(\sigma_{23}\) obstruction is closed.
\end{theorem}

\begin{proof}
Fix \(x\in S_R\).

For \(\sigma_{21}\), the right word is \([1,2,2]\).  Since \(x\in R1\), the first step is
\[
\sigma_{21}(x)=\stepd{1}(x).
\]
By the previous lemma this target lies in \(R2\), and \(R2\) is routed by direction \(2\).  Therefore
\[
\sigma_{21}^2(x)=\stepd{2}\stepd{1}(x).
\]

For \(\sigma_{22}\), the right word is \([2,1,2]\).  The first two right-side directions are therefore reversed:
\[
\sigma_{22}(x)=\stepd{2}(x),\qquad \sigma_{22}^2(x)=\stepd{1}\stepd{2}(x).
\]
Because the surfaced common module defines each \(\stepd{d}\) as coordinatewise \(+1\) in a single axis, the
steps commute, so
\[
\sigma_{22}^2(x)=\stepd{1}\stepd{2}(x)=\stepd{2}\stepd{1}(x)=\sigma_{21}^2(x).
\]

For \(\sigma_{23}\), the right word is \([2,2,1]\), so
\[
\sigma_{23}(x)=\stepd{2}(x),\qquad \sigma_{23}^2(x)=\stepd{2}^2(x).
\]
Now every source in \(S_R\) has \((q,w)=(m-1,0)\).  Hence
\[
\sigma_{21}^2(x): (q,w)=(0,1),
\qquad
\sigma_{23}^2(x): (q,w)=(m-1,2).
\]
These are distinct for every \(m>2\).  Therefore
\[
\sigma_{21}^2(x)=\sigma_{22}^2(x)\neq \sigma_{23}^2(x),
\]
as required.
\end{proof}

\begin{remark}[What changed relative to the earlier frontier]
Before the bundle was surfaced, the live burden was to identify the hidden immediate rule or to certify enough
bounded locality to transfer the checked range to all odd \(m\).  After surfacing the builder, the relevant rule is
no longer hidden: the second-step comparison is an explicit consequence of \texttt{\_prepare\_m},
\texttt{\_build\_candidate}, and the commuting coordinate steps from \texttt{precompute\_m}.  The residual task is
not \(\sigma_{23}\) itself, but only the manuscript-side translation of this surfaced-script proof into the preferred
notation of the main text.
\end{remark}

\paragraph{Compute sanity check.}
The companion rerun script \texttt{d5\_219\_sigma23\_bundle\_exact\_checker.py} checks the surfaced theorem on the explicit
range
\[
m=3,5,7,9,11,13,15,17,19,
\]
and finds exactly one rooted microcase at each modulus, source count \(m(m-1)\), neighbor pattern
\([B,R2,B,B,B]\), relative immediate displacement \((0,1,0,0,0)\), relative actual second displacement
\((0,1,1,0,0)\), and persistent avoidance of \(\sigma_{23}\).

\end{document}
